% Auto-generated by scripts/run_spfpp_ablation.py
\subsection{Point-Based Semantic Matching (No Semantic Walls)}
To ablate semantic walls, we model poles and trunks as individual map objects and score each observation ray against class-consistent point hypotheses.

Given particle $i$ and observation ray $j$ with class $c_j$, observed range $r_j$, and local bearing $\alpha_j$, we transform all map points of class $c\in\{\text{pole},\text{trunk}\}$ into the particle local frame, obtaining $(r_{ik},\alpha_{ik})$.

Candidates satisfy:
\begin{equation}
|\Delta\alpha_{ijk}| \le \alpha_{\text{gate}},\qquad
|\Delta r_{ijk}| \le r_{\text{gate}},\qquad
0 < r_{ik} \le r_{\max},
\end{equation}
with $\Delta\alpha_{ijk}=\mathrm{wrap}(\alpha_{ik}-\alpha_j)$ and $\Delta r_{ijk}=r_{ik}-r_j$.

If class-consistent candidates exist, we use the best (maximum log-likelihood):
\begin{equation}
s_{ij}=\max_k\left[-\frac{1}{2}\left(\frac{\Delta\alpha_{ijk}^2}{\sigma_\alpha^2}+\frac{\Delta r_{ijk}^2}{\sigma_r^2}\right)\right].
\end{equation}
If only opposite-class candidates exist, a wrong-hit penalty is used; if no candidate exists, a miss penalty is used.

For background rays, if any map point falls along the ray before the observed background range, we treat the ray as a hit-like free-space consistency event; otherwise a miss penalty is applied.

The per-particle semantic score is averaged across rays and fused with GPS and corridor terms using the same normalized fusion strategy used in SPF++.
